\documentclass[11pt,letterpaper]{article}
\usepackage{../../latex_common/estilo_monitoria}

\begin{document}

\portadamodulo{00}{Algoritmos para Tontos (Pensar Antes de Teclar)}{2.5 horas}{%
  Lógica cotidiana, modelo Entrada--Proceso--Salida (E-P-S), Método de los 5 Pasos y prueba de escritorio en papel%
}{%
  \begin{enumerate}
    \item Perderle el miedo a la palabra \textbf{algoritmo} y entender exactamente cómo ``piensa'' una computadora.
    \item Leer cualquier enunciado de clase o parcial y separar sus piezas con el modelo \textbf{E-P-S} y el diccionario de palabras clave.
    \item Aplicar el \textbf{Método de los 5 Pasos del Monitor} para resolver problemas cotidianos y matemáticos sin quedarse en blanco frente al teclado.
    \item Ejecutar una \textbf{prueba de escritorio en papel} para cazar errores de lógica antes de escribir una sola línea de código.
  \end{enumerate}
}

\tableofcontents
\vspace{8pt}

\section{Mapa Conceptual del Módulo}

Antes de entrar en materia, mira la ruta completa que vas a recorrer hoy. El secreto de los mejores programadores no está en teclear rápido, sino en seguir este mapa de pensamiento ordenado antes de tocar el computador:

\begin{center}
\begin{tikzpicture}[
  node distance=0.55cm,
  caja/.style={rectangle, rounded corners=4pt, draw=AzulMarino, fill=AzulClaro, thick,
               text width=0.78\linewidth, align=center, font=\sffamily\small, inner sep=7pt},
  flecha/.style={-{Stealth[length=2.5mm]}, thick, color=Turquesa}
]
  \node[caja, fill=AzulMarino, text=white] (n0) {\textbf{ENUNCIADO DEL PROBLEMA} (En español, a veces largo o confuso)};
  \node[caja, below=of n0] (n1) {\textbf{Paso 1: Leer con lupa y subrayar} (datos variables, constantes y reglas ocultas)};
  \node[caja, below=of n1, fill=TurquesaClaro, draw=Turquesa] (n2) {\textbf{Paso 2: Modelo E-P-S (La Licuadora)} \\ \texttt{ENTRADAS} $\longrightarrow$ \texttt{PROCESO} $\longrightarrow$ \texttt{SALIDAS}};
  \node[caja, below=of n2] (n3) {\textbf{Paso 3: Elegir la estructura lógica} (Secuencia $\mid$ Decisión Sí/No $\mid$ Ciclo)};
  \node[caja, below=of n3, fill=NaranjaClaro, draw=NaranjaCalido] (n4) {\textbf{Paso 4: Prueba de escritorio en papel} (Simular con datos normales y casos tramposos)};
  \node[caja, below=of n4, fill=VerdeClaro, draw=VerdeExito] (n5) {\textbf{Paso 5: Traducción futura} $\longrightarrow$ DFD (M01) $\longrightarrow$ PSeInt (M02) $\longrightarrow$ C++ (M02+)};

  \draw[flecha] (n0) -- (n1);
  \draw[flecha] (n1) -- (n2);
  \draw[flecha] (n2) -- (n3);
  \draw[flecha] (n3) -- (n4);
  \draw[flecha] (n4) -- (n5);
\end{tikzpicture}
\end{center}

\section{Explicado ``Para Tontos'': ¿Qué es un Algoritmo?}

\imagenmodulo[0.86\linewidth]{ilus_01_algoritmo_receta.jpg}{Un algoritmo es simplemente una receta ordenada paso a paso para darle instrucciones exactas a la computadora.}

\begin{paratontos}{¿De qué va esto sin palabras raras?}
Si alguien te dice \textit{``diseña un algoritmo eficiente''}, suena a película de ciencia ficción o a fórmula de la NASA. Respira tranquilo: \textbf{tú ya usas algoritmos todos los días desde que eras niño}.

Un \textbf{algoritmo} es simplemente una \textbf{lista ordenada, finita y sin ambigüedades de pasos} para lograr un objetivo o resolver un problema.
\begin{itemize}
  \item \textbf{Ordenada:} El orden importa. No puedes ponerte los zapatos antes que las medias.
  \item \textbf{Finita:} Tiene que terminar en algún momento. Si un paso te manda en círculos para siempre, nunca acabas.
  \item \textbf{Sin ambigüedades:} Cero adivinanzas. No le digas a un robot \textit{``échale sal al gusto''}, porque el robot te preguntará: \textit{``¿qué es gusto? ¿1 gramo o medio kilo?''}.
\end{itemize}
\end{paratontos}

\subsection{El secreto mejor guardado: la computadora es un ``Genio Literal''}

Imagina que tienes un asistente súper veloz capaz de hacer un millón de sumas por segundo, pero \textbf{tiene cero sentido común}. Si le dices:
\textit{``Unta mantequilla en el pan''},
y no le dijiste primero \textit{``abre la nevera, saca la mantequilla, toma un cuchillo, destapa el recipiente...''}, el genio literal aplastará el pan contra el frasco cerrado de mantequilla y te mirará orgulloso diciendo: \textit{``¡Misión cumplida!''}.

\begin{consejo}{La regla mental del monitor}
Programar \textbf{NO} es ser un genio de las matemáticas; programar es \textbf{aprender a darle instrucciones tan claras y ordenadas a un robot literal, que sea imposible que se equivoque}.
\end{consejo}

\section{Cinco Analogías Completas de la Vida Real}

Para que el concepto quede grabado con plastilina, analicemos cinco situaciones cotidianas exactamente como las piensa un programador.

\subsection{Analogía 1: Preparar un café con leche por la mañana (Secuencia estricta)}
\begin{itemize}
  \item \textbf{Objetivo:} Obtener una taza de café con leche caliente.
  \item \textbf{Entradas (ingredientes):} Agua, café instantáneo, leche, azúcar, taza, cuchara, hervidor.
  \item \textbf{Pasos del algoritmo:}
    \begin{enumerate}
      \item Inicio.
      \item Poner a calentar el agua en el hervidor hasta que hierva.
      \item Colocar la taza limpia sobre la mesa.
      \item Agregar 1 cucharadita de café en la taza.
      \item Verter el agua caliente hasta llenar $3/4$ de la taza.
      \item Agregar leche hasta completar la taza y 1 cucharadita de azúcar.
      \item Revolver con la cuchara durante 10 segundos.
      \item Entregar la taza lista (\textbf{Salida}) y Fin.
    \end{enumerate}
  \item \textbf{¿Qué pasa si cambias el orden?} Si viertes el agua caliente antes de colocar la taza (intercambias pasos 3 y 5), terminas con la mesa quemada y la taza vacía. En programación, el orden de las líneas lo cambia todo.
\end{itemize}

\subsection{Analogía 2: Pedir una hamburguesa por una app de domicilios (Decisiones)}
\begin{itemize}
  \item \textbf{Objetivo:} Recibir tu comida en casa pagando correctamente.
  \item \textbf{Entradas:} Dirección de entrega, restaurante elegido, producto, método de pago y saldo disponible.
  \item \textbf{Pasos del algoritmo:}
    \begin{enumerate}
      \item Inicio $\rightarrow$ Abrir la app e ingresar dirección.
      \item \textbf{Decisión 1:} ¿El restaurante está abierto? Si \textbf{No}, mostrar ``Restaurante cerrado'' y terminar; si \textbf{Sí}, continuar.
      \item Agregar la hamburguesa al carrito y elegir pago (tarjeta o efectivo).
      \item \textbf{Decisión 2:} ¿Pagas con tarjeta y tienes saldo suficiente? Si \textbf{Sí}, confirmar pedido y descontar dinero; si \textbf{No}, pedir otro método o cancelar.
      \item Esperar al domiciliario, recibir pedido (\textbf{Salida}) y Fin.
    \end{enumerate}
  \item \textbf{Lección clave:} Los algoritmos reales no siempre van en línea recta; toman caminos distintos según una condición de \textbf{Sí} o \textbf{No}.
\end{itemize}

\subsection{Analogía 3: Retirar dinero en un cajero automático (Seguridad y repetición)}
\begin{itemize}
  \item \textbf{Objetivo:} Sacar efectivo sin entregar dinero de más ni permitir fraudes.
  \item \textbf{Entradas:} Tarjeta débito, clave PIN, monto a retirar, saldo de la cuenta.
  \item \textbf{Pasos del algoritmo:}
    \begin{enumerate}
      \item Inicio $\rightarrow$ Insertar tarjeta e iniciar un contador de intentos en \codigo{0}.
      \item Pedir al usuario que digite su PIN de 4 dígitos.
      \item Si el PIN es correcto, pasar al paso 5. Si es incorrecto, sumar \codigo{1} al contador.
      \item Si el contador es menor que \codigo{3}, volver al paso 2; si llegó a \codigo{3}, retener tarjeta y terminar por seguridad.
      \item Leer el monto solicitado. \textbf{Decisión:} ¿El monto es menor o igual al saldo \textbf{Y} múltiplo de \$10.000? Si \textbf{Sí}, restar del saldo y entregar billetes y recibo (\textbf{Salidas}); si \textbf{No}, mostrar error.
      \item Devolver la tarjeta y Fin.
    \end{enumerate}
  \item \textbf{Lección clave:} Aquí ves \textbf{las 3 grandes piezas} de la programación juntas: \textbf{secuencia}, \textbf{decisión} y \textbf{ciclo/repetición} (intentar el PIN hasta 3 veces).
\end{itemize}

\subsection{Analogía 4: Armar un mueble o figura LEGO con instructivo (Orden y variables)}
\begin{itemize}
  \item \textbf{Objetivo:} Construir la silla exactamente igual a la foto de la caja.
  \item \textbf{Entradas:} Bolsas de piezas etiquetadas (A, B, C), tornillos, destornillador y manual.
  \item \textbf{Pasos:} Contar que estén todas las piezas $\rightarrow$ Seguir el bloque 1 antes del bloque 2 $\rightarrow$ Repetir el ajuste para las 4 patas (ciclo de 4 veces) $\rightarrow$ Verificar estabilidad.
  \item \textbf{Lección clave:} El manual numera las bolsas para no usar una pieza que todavía no sacaste. En programación, si multiplicas una variable antes de leer su valor, es como atornillar una pata invisible.
\end{itemize}

\subsection{Analogía 5: Hacer la maleta para un viaje según el clima (Ciclos + Decisiones)}
\begin{itemize}
  \item \textbf{Objetivo:} Empacar la ropa exacta para 4 días.
  \item \textbf{Entradas:} Días del viaje (\codigo{4}), pronóstico del clima (\texttt{"lluvioso"} o \texttt{"soleado"}).
  \item \textbf{Pasos:} Abrir maleta $\rightarrow$ \textbf{Repetir 4 veces}: meter 1 camiseta y 1 par de medias $\rightarrow$ \textbf{Decisión}: ¿va a llover? Si \textbf{Sí}, meter paraguas y chaqueta impermeable; si \textbf{No}, meter gorra y bloqueador $\rightarrow$ Cerrar maleta (\textbf{Salida}).
  \item \textbf{Lección clave:} Un ciclo (\textit{``repetir N veces''}) evita escribir cuatro veces la misma orden. Si un proceso ocurre 100 veces, usas un ciclo.
\end{itemize}

\section{Cómo Analizar Problemas Desde Cero: El Modelo E-P-S}

Cuando un profesor te entrega un enunciado largo en un parcial, el error número uno del estudiante novato es abrir el computador de inmediato y empezar a teclear sin entender qué le pidieron. \textbf{Antes de pensar en código, piensa como una licuadora:}

\begin{center}
\begin{tikzpicture}[
  bloque/.style={rectangle, rounded corners=5pt, thick, minimum height=1.35cm, text width=4.2cm, align=center, font=\sffamily\small},
  flecha/.style={-{Stealth[length=3mm]}, ultra thick, color=AzulMarino}
]
  \node[bloque, draw=Turquesa, fill=TurquesaClaro] (e) {\textbf{\color{Turquesa}1. ENTRADAS (E)}\\Datos que vienen de afuera\\(teclado / ingredientes)};
  \node[bloque, right=1.0cm of e, draw=NaranjaCalido, fill=NaranjaClaro] (p) {\textbf{\color{NaranjaCalido}2. PROCESO (P)}\\Fórmulas, cálculos y\\decisiones internas};
  \node[bloque, right=1.0cm of p, draw=VerdeExito, fill=VerdeClaro] (s) {\textbf{\color{VerdeExito}3. SALIDAS (S)}\\Resultados finales que\\ve el usuario en pantalla};
  \draw[flecha] (e) -- (p);
  \draw[flecha] (p) -- (s);
\end{tikzpicture}
\end{center}

\begin{consejo}{Lee de atrás hacia adelante si te bloqueas}
Empieza identificando \textbf{la SALIDA} (\textit{¿qué me piden mostrar al final?}). Luego pregúntate: \textit{``¿Qué fórmula o condición produce esa salida?''} (ese es tu \textbf{PROCESO}). Por último pregúntate: \textit{``¿Qué datos variables le faltan a esa fórmula?''} (esas son tus \textbf{ENTRADAS}).
\end{consejo}

\subsection{Diccionario Detective: Palabras Clave en los Enunciados}

Cuando leas un enunciado de taller o parcial, busca estas pistas como un detective:

\begin{center}
\small
\renewcommand{\arraystretch}{1.3}
\rowcolors{2}{GrisSuave}{white}
\begin{tabularx}{\linewidth}{>{\leavevmode\raggedright\arraybackslash}p{5.5cm} >{\leavevmode\bfseries\sffamily\color{AzulMarino}\raggedright\arraybackslash}p{3.0cm} >{\leavevmode\raggedright\arraybackslash}X}
  \rowcolor{AzulMarino}
  \textcolor{white}{\textbf{Si el enunciado dice...}} &
  \textcolor{white}{\textbf{Suele ser...}} &
  \textcolor{white}{\textbf{Qué significa en tu diseño}} \\
  \textit{``Leer'', ``Ingresar'', ``Pedir al usuario'', ``Digitar'', ``Capturar'', ``Dado un número...''} & ENTRADA & Creas una variable (cajita) y pides ese valor al usuario por teclado. \\
  \textit{``Calcular'', ``Determinar'', ``Sumar'', ``Promediar'', ``Aplicar descuento'', ``Convertir'', ``Verificar si...''} & PROCESO & Aquí van tus operaciones aritméticas o tus decisiones (\textit{Si... entonces}). \\
  \textit{``Mostrar'', ``Imprimir'', ``Escribir en pantalla'', ``Reportar'', ``Desplegar'', ``Indicar si...''} & SALIDA & Es el mensaje o valor numérico final que aparece en la consola. \\
  \textit{``El IVA es del 19\%'', ``1 hora = 60 min'', ``Se aprueba con 3.0 o más''} & CONSTANTE / REGLA FIJA & \textbf{¡NO se pide por teclado!} Ya te lo dio fijo el enunciado para tu proceso. \\
\bottomrule
\end{tabularx}
\end{center}

\begin{peligro}{Pedir constantes por teclado}
Muchos estudiantes leen: \textit{``Calcule el total a pagar aplicando un IVA del 19\%''} y le piden al usuario que digite el \codigo{19\%}. \textbf{¡Incorrecto!} Si el problema ya fija el 19\%, ese dato es una constante del proceso, no una entrada. Solo pides por teclado lo que cambia entre un usuario y otro (el precio del producto).
\end{peligro}

\section{El Método de los 5 Pasos del Monitor}

\imagenmodulo[0.95\linewidth]{met_01_cinco_pasos_analisis.png}{Infografía del Método de los 5 Pasos del Monitor: analiza y prueba en papel antes de abrir el editor.}

Cada vez que tengas un ejercicio enfrente —desde hoy hasta tu examen final— aplica esta receta en orden estricto:

\begin{enumerate}
  \item \textbf{Paso 1 -- Leer con lupa y subrayar:} Lee dos veces. Subraya los datos que cambian (entradas), las constantes fijas, las reglas (\textit{``si ocurre...''}) y lo que te piden entregar.
  \item \textbf{Paso 2 -- Armar la Tabla E-P-S y bautizar variables:} Anota Entradas, Proceso y Salidas, y pon nombres claros (\codigo{precioProducto}, \codigo{notaFinal}; jamás \codigo{x1} o \codigo{cosa}).
  \item \textbf{Paso 3 -- Elegir la estructura lógica:} ¿Pasos en fila? $\rightarrow$ \textbf{Secuencia}. ¿Hay un ``si pasa X''? $\rightarrow$ \textbf{Decisión}. ¿Hay que repetir algo? $\rightarrow$ \textbf{Ciclo}.
  \item \textbf{Paso 4 -- Hacer la Prueba de Escritorio en papel:} Inventa datos fáciles, simula cada línea en una tabla y prueba también un \textbf{caso tramposo} (un \codigo{0}, un valor exacto en el límite o un negativo).
  \item \textbf{Paso 5 -- Traducir a DFD, PSeInt o C++:} El código es solo la traducción final; el 80\% del éxito ocurre en los pasos 1 al 4.
\end{enumerate}

\imagenmodulo[0.86\linewidth]{ilus_03_puente_dfd_pseint_cpp.jpg}{Cuando los pasos 1 al 4 están claros en papel, traducir a DFD, PSeInt y C++ se vuelve un puente natural.}

\section{Método de los 5 Pasos Aplicado a 4 Problemas Completos}

\subsection{Problema 1 (Cotidiano): Repartir la cuenta de una pizza entre amigos}
\begin{analisis}{Enunciado -- La pizza del viernes}
Un grupo de amigos pide una pizza a domicilio. El programa debe recibir el valor de la pizza, el valor del domicilio, el porcentaje de propina (por ejemplo \codigo{10} para 10\%) y la cantidad de amigos que pagan. Calcula el total de la cuenta (aplicando la propina sobre pizza + domicilio) y muestra cuánto debe poner cada amigo en partes iguales.
\end{analisis}

\begin{itemize}
  \item \textbf{Paso 1 (Subrayar):} Entradas variables: valor pizza, valor domicilio, porcentaje propina, cantidad de amigos. Salidas: total de la cuenta y cuota por amigo.
  \item \textbf{Paso 2 (Tabla E-P-S):}
    \begin{itemize}
      \item \textbf{Entradas:} \codigo{valorPizza}, \codigo{valorDomicilio}, \codigo{porcPropina}, \codigo{numAmigos}.
      \item \textbf{Proceso:} \texttt{subtotal = valorPizza + valorDomicilio};\\
            \texttt{propina = subtotal * (porcPropina / 100.0)}; \texttt{totalCuenta = subtotal + propina};\\
            \texttt{cuotaAmigo = totalCuenta / numAmigos} (suponiendo \codigo{numAmigos > 0}).
      \item \textbf{Salidas:} \codigo{totalCuenta}, \codigo{cuotaAmigo}.
    \end{itemize}
  \item \textbf{Paso 3 (Estructura):} \textbf{Secuencial} (un cálculo tras otro).
\end{itemize}

\begin{pruebaescritorio}{Problema 1 -- Datos: pizza=45000, domicilio=5000, propina=10\%, amigos=5}
\begin{center}
\footnotesize
\setlength{\tabcolsep}{4pt}
\renewcommand{\arraystretch}{1.25}
\begin{tabular}{c l c c c c c c c c}
  \toprule
  \textbf{Paso} & \textbf{Acción} & \textbf{vPizza} & \textbf{vDom} & \textbf{\%Prop} & \textbf{amigos} & \textbf{subtotal} & \textbf{propina} & \textbf{total} & \textbf{cuota} \\
  \midrule
  1 & Leer entradas & 45000 & 5000 & 10 & 5 & -- & -- & -- & -- \\
  2 & Calcular subtotal & 45000 & 5000 & 10 & 5 & \textbf{50000} & -- & -- & -- \\
  3 & Calcular propina & 45000 & 5000 & 10 & 5 & 50000 & \textbf{5000} & -- & -- \\
  4 & Calcular totalCuenta & 45000 & 5000 & 10 & 5 & 50000 & 5000 & \textbf{55000} & -- \\
  5 & Calcular cuotaAmigo & 45000 & 5000 & 10 & 5 & 50000 & 5000 & 55000 & \textbf{11000} \\
  \bottomrule
\end{tabular}
\end{center}
\textbf{Salida en pantalla:} Total de la cuenta: \$55000 $\mid$ Cuota por amigo: \$11000.
\end{pruebaescritorio}

\textbf{Paso 5 (Plan en español):} Leer las 4 entradas $\rightarrow$ calcular \codigo{subtotal}, \codigo{propina}, \codigo{totalCuenta} y\\
\codigo{cuotaAmigo} $\rightarrow$ mostrar \codigo{totalCuenta} y \codigo{cuotaAmigo}.

\subsection{Problema 2 (Cotidiano con Decisión): ¿Me alcanza para ir al cine?}
\begin{analisis}{Enunciado -- Salida al cine}
La boleta de cine cuesta \$18.000 fijos, pero si pides un combo debes sumar su precio. Lee cuánto dinero tienes ahorrado (\codigo{dineroDisponible}) y el precio del combo (\codigo{precioCombo}). Calcula el costo total y dime si te alcanza el dinero (indicando cuánto te sobra o cuánto te falta).
\end{analisis}

\begin{itemize}
  \item \textbf{Paso 1:} Constante oculta: \texttt{PRECIO\_ENTRADA = 18000} (no se pide por teclado).
  \item \textbf{Paso 2 (E-P-S):} Entradas: \codigo{dineroDisponible}, \codigo{precioCombo}.\\
        \textbf{Proceso:} \texttt{costoTotal = 18000 + precioCombo}.\\
        Si \texttt{dineroDisponible >= costoTotal}, \texttt{sobrante = dineroDisponible - costoTotal};\\
        si no, \texttt{faltante = costoTotal - dineroDisponible}.\\
        \textbf{Salidas:} \codigo{costoTotal} y mensaje con sobrante/faltante.
  \item \textbf{Paso 3 (Estructura):} \textbf{Secuencia + Decisión Simple} (observa que usamos \codigo{>=}: si tienes el dinero exacto, sí te alcanza y te sobra \$0).
\end{itemize}

\begin{pruebaescritorio}{Problema 2 -- Probamos ambos caminos}
\begin{itemize}
  \item \textbf{Caso A (Alcanza):} \texttt{dineroDisponible = 50000}, \texttt{precioCombo = 24000}.\\
        \texttt{costoTotal = 18000 + 24000 = 42000}. ¿\texttt{50000 >= 42000}? \textbf{Sí} $\rightarrow$ \texttt{sobrante = 8000}.\\
        \textit{Pantalla:} \texttt{Costo: \$42000. ¡Sí te alcanza! Te sobran \$8000.}
  \item \textbf{Caso B (No alcanza):} \texttt{dineroDisponible = 40000}, \texttt{precioCombo = 25000}.\\
        \texttt{costoTotal = 43000}. ¿\texttt{40000 >= 43000}? \textbf{No} $\rightarrow$ \texttt{faltante = 3000}.\\
        \textit{Pantalla:} \texttt{Costo: \$43000. No te alcanza. Te faltan \$3000.}
\end{itemize}
\end{pruebaescritorio}

\subsection{Problema 3 (Académico): Nota definitiva ponderada y aprobación}
\begin{analisis}{Enunciado -- Tres cortes universitarios}
El Primer Corte vale 30\%, el Segundo Corte vale 30\% y el Tercer Corte vale 40\%. Se aprueba con \codigo{3.0} o más. Lee las tres notas, calcula la definitiva y muestra si el estudiante fue \texttt{"APROBADO"} o \texttt{"REPROBADO"}.
\end{analisis}

\begin{itemize}
  \item \textbf{Paso 2 (E-P-S):} Entradas: \codigo{nota1}, \codigo{nota2}, \codigo{nota3}.\\
        \textbf{Proceso:} \texttt{definitiva = nota1*0.30 + nota2*0.30 + nota3*0.40};\\
        evaluar \codigo{definitiva >= 3.0}.\\
        \textbf{Salidas:} \codigo{definitiva} y estado.
  \item \textbf{Paso 3:} \textbf{Secuencia + Decisión Simple} con \codigo{>= 3.0}.
\end{itemize}

\begin{pruebaescritorio}{Problema 3 -- Incluye caso borde exacto 3.0}
\begin{itemize}
  \item \textbf{Caso A (Normal):} \codigo{nota1 = 2.0, nota2 = 3.0, nota3 = 4.0} $\rightarrow$ \codigo{definitiva = 0.60 + 0.90 + 1.60 = 3.10} $\rightarrow$ \codigo{3.10 >= 3.0} (V) $\rightarrow$ \texttt{APROBADO}.
  \item \textbf{Caso B (Caso límite):} \codigo{3.0, 3.0, 3.0} $\rightarrow$ \codigo{definitiva = 3.00} $\rightarrow$ \codigo{3.00 >= 3.0} (V) $\rightarrow$ \texttt{APROBADO}. (Si hubieras puesto \codigo{> 3.0}, habría reprobado injustamente.)
  \item \textbf{Caso C (Reprueba):} \codigo{2.5, 2.5, 3.2} $\rightarrow$ \codigo{definitiva = 0.75 + 0.75 + 1.28 = 2.78} $\rightarrow$ \codigo{2.78 >= 3.0} (F) $\rightarrow$ \texttt{REPROBADO}.
\end{itemize}
\end{pruebaescritorio}

\subsection{Problema 4 (Física básica): Tiempo de viaje en carretera}
\begin{analisis}{Enunciado -- Bus intermunicipal}
Conocida la distancia en kilómetros (\codigo{distanciaKm}) y la velocidad promedio en km/h (\codigo{velocidadKmH}), calcula cuánto tarda el viaje en \textbf{horas} y en \textbf{minutos}. Valida que la velocidad sea mayor a cero.
\end{analisis}

\begin{itemize}
  \item \textbf{Paso 2 (E-P-S):} Entradas: \codigo{distanciaKm}, \codigo{velocidadKmH}. Proceso: si \codigo{velocidadKmH > 0}, calcular \codigo{tiempoHoras = distanciaKm / velocidadKmH} y \codigo{tiempoMinutos = tiempoHoras * 60.0}; si no, reportar error. Salidas: \codigo{tiempoHoras} y \codigo{tiempoMinutos} (o mensaje de error).
  \item \textbf{Paso 3:} \textbf{Decisión Simple de validación + Secuencia}.
\end{itemize}

\begin{pruebaescritorio}{Problema 4 -- Caso normal y caso tramposo velocidad = 0}
\begin{itemize}
  \item \textbf{Caso A:} \codigo{distanciaKm = 150}, \codigo{velocidadKmH = 60}. Como \codigo{60 > 0} (V) $\rightarrow$ \codigo{tiempoHoras = 150 / 60 = 2.5} h; \codigo{tiempoMinutos = 2.5 * 60 = 150} min.
  \item \textbf{Caso B (Tramposo):} \codigo{distanciaKm = 150}, \codigo{velocidadKmH = 0}. Como \codigo{0 > 0} (F) $\rightarrow$ \textit{Pantalla:} \texttt{Error: la velocidad debe ser mayor a cero.} ¡Evitamos explotar por división entre cero!
\end{itemize}
\end{pruebaescritorio}

\section{Consejos de Monitoría y Errores Clásicos}

\begin{consejo}{Hábitos que garantizan pasar el primer parcial}
\begin{enumerate}
  \item \textbf{No abras el editor los primeros 5 minutos:} Si no puedes explicar el problema en papel con el Método de los 5 Pasos, todavía no estás listo para programar.
  \item \textbf{Bautiza bien tus cajas:} Usa \codigo{precioNeto} o \codigo{edadUsuario}, jamás \codigo{x1} o \codigo{a}.
  \item \textbf{Pregunta siempre por el caso cero y el empate:} Prueba tus fórmulas con \codigo{0}, con números exactamente en el límite y con valores iguales.
\end{enumerate}
\end{consejo}

\begin{peligro}{Los 5 errores clásicos de principiante (y spoiler de C++)}
\begin{enumerate}
  \item \textbf{Confundir asignación con igualdad:} En matemáticas \codigo{x = x + 1} no tiene sentido; en programación \codigo{x <- x + 1} significa \textit{``toma el valor actual de x, súmale 1 y guárdalo encima en x''}.
  \item \textbf{Olvidar paréntesis:} Escribir \codigo{a + b + c / 3} divide primero \codigo{c / 3} y daña el promedio. Debe ser \codigo{(a + b + c) / 3}.
  \item \textbf{Leer por teclado la salida:} Si te piden calcular el área de un triángulo, \textbf{no} hagas \codigo{Leer area}; pide \codigo{base} y \codigo{altura}.
  \item \textbf{Saltarse la prueba de escritorio:} Creer que ``a ojo'' funciona y descubrir en el parcial que multiplicaste donde tocaba dividir.
  \item \textbf{Tenerle miedo a los errores rojos de C++:} Cuando llegues al Módulo 02 y el compilador te diga en inglés cosas como \texttt{expected ';' before ...} o \texttt{was not declared in this scope}, recuerda que solo significa \textit{``olvidaste un punto y coma en la línea anterior''} o \textit{``usaste una caja sin crearla primero''}. ¡En el M02 tendrás el diccionario completo!
\end{enumerate}
\end{peligro}

\section{Ejercicios Propuestos (Lógica y Análisis)}

Resuelve estos \textbf{8 ejercicios en papel} aplicando el Método de los 5 Pasos completo antes de consultar el solucionario al final del módulo.

\begin{ejercicio}{1 -- Receta escalable de panqueques}
Una receta para \textbf{4 personas} requiere 2 huevos, 200 g de harina y 300 ml de leche. Diseña el algoritmo que reciba cuántas personas van a desayunar y calcule cuántos huevos, gramos de harina y ml de leche necesitas.
\end{ejercicio}

\begin{ejercicio}{2 -- El estudiante distraído (Detective E-P-S)}
Para el enunciado \textit{``Calcular el precio final de un pantalón sabiendo que tiene un descuento fijo del 20\% sobre su precio original''}, un estudiante puso como entradas: \codigo{precioOriginal}, \codigo{descuento = 20}, \codigo{precioFinal}. Identifica sus errores y escribe la tabla E-P-S y el algoritmo correctos.
\end{ejercicio}

\begin{ejercicio}{3 -- Conversión de horas}
Dado un número entero de horas (ejemplo \codigo{5}), calcula a cuántos minutos y a cuántos segundos equivale.
\end{ejercicio}

\begin{ejercicio}{4 -- Cambio o faltante en la tienda}
Un cliente compra un producto de precio \codigo{P} y paga con un billete \codigo{B}. Diseña el análisis y algoritmo para saber si alcanza y cuánto cambio entregarle, o cuánto dinero le falta.
\end{ejercicio}

\begin{ejercicio}{5 -- Intercambio de vasos (Jugo y Agua)}
Tienes \codigo{vasoA = "Jugo"} y \codigo{vasoB = "Agua"}. Escribe el algoritmo paso a paso para intercambiar sus contenidos sin perder ninguno (pista: usa una tercera variable auxiliar) y haz su prueba de escritorio.
\end{ejercicio}

\begin{ejercicio}{6 -- Edad y mayoría de edad}
Dado el año de nacimiento de una persona y el año actual, calcula su edad aproximada e indica si es mayor de edad (\codigo{>= 18} años) o menor de edad.
\end{ejercicio}

\begin{ejercicio}{7 -- Cazar el error en la prueba de escritorio}
Un compañero escribió: \texttt{1. Leer a, b} $\rightarrow$ \texttt{2. promedio <- a + b / 2} $\rightarrow$ \texttt{3. Escribir promedio}. Realiza la prueba de escritorio con \codigo{a = 10} y \codigo{b = 20}, encuentra por qué falla y corrígelo.
\end{ejercicio}

\begin{ejercicio}{8 -- Consumo de gasolina en moto}
Conoces la distancia de un viaje en km, cuántos kilómetros rinde la moto por galón y el precio de un galón. Calcula cuántos galones consumirás y cuánto dinero gastarás en gasolina.
\end{ejercicio}

\section{Solucionario Explicado Paso a Paso}

\begin{solucion}{Ejercicios 1 al 4}
\textbf{Ejercicio 1 (Panqueques):}
\begin{itemize}
  \item \textbf{E-P-S:} Entrada: \codigo{personas}. Proceso: \codigo{factor = personas / 4.0}, \codigo{huevos = 2 * factor}, \codigo{harinaGr = 200 * factor}, \codigo{lecheMl = 300 * factor}. Salidas: \codigo{huevos}, \codigo{harinaGr}, \codigo{lecheMl}.
  \item \textbf{Prueba (\codigo{personas = 6}):} \codigo{factor = 1.5} $\rightarrow$ \codigo{huevos = 3.0}, \codigo{harinaGr = 300}, \codigo{lecheMl = 450}.
\end{itemize}

\textbf{Ejercicio 2 (Estudiante distraído):}
\begin{itemize}
  \item \textbf{Errores:} (1) Pidió \codigo{descuento = 20} como entrada cuando el 20\% es constante fija; (2) puso \codigo{precioFinal} como entrada cuando es la \textbf{salida}.
  \item \textbf{Corrección:} Entrada: \codigo{precioOriginal}. Proceso: \codigo{descuento = precioOriginal * 0.20}, \codigo{precioFinal = precioOriginal - descuento}. Prueba con \codigo{100000}: \codigo{descuento = 20000}, \codigo{precioFinal = 80000}.
\end{itemize}

\textbf{Ejercicio 3 (Horas a minutos y segundos):}
\begin{itemize}
  \item \textbf{E-P-S:} Entrada: \codigo{horas}. Proceso: \codigo{minutos = horas * 60}, \codigo{segundos = minutos * 60}. Salidas: \codigo{minutos}, \codigo{segundos}.
  \item \textbf{Prueba (\codigo{horas = 5}):} \codigo{minutos = 300}, \codigo{segundos = 18000}.
\end{itemize}

\textbf{Ejercicio 4 (Cambio en tienda):}
\begin{itemize}
  \item \textbf{E-P-S:} Entradas: \codigo{precioProducto}, \codigo{pagoCliente}. Proceso: si \codigo{pagoCliente >= precioProducto}, \codigo{cambio = pagoCliente - precioProducto}; si no, \codigo{faltante = precioProducto - pagoCliente}.
  \item \textbf{Prueba:} Con \codigo{P = 12500, B = 20000} $\rightarrow$ \codigo{cambio = 7500}. Con \codigo{P = 12500, B = 10000} $\rightarrow$ \codigo{faltante = 2500}.
\end{itemize}
\end{solucion}

\begin{solucion}{Ejercicios 5 al 8}
\textbf{Ejercicio 5 (Intercambio de vasos):} Si haces \codigo{vasoA <- vasoB} directo, pierdes el jugo. Usa \codigo{vasoExtra}:
\begin{enumerate}
  \item \codigo{vasoExtra <- vasoA} \textit{(guarda ``Jugo'')}
  \item \codigo{vasoA <- vasoB} \textit{(ahora vasoA tiene ``Agua'')}
  \item \codigo{vasoB <- vasoExtra} \textit{(ahora vasoB tiene ``Jugo'')}
\end{enumerate}

\textbf{Ejercicio 6 (Edad y mayoría de edad):}
\begin{itemize}
  \item \textbf{E-P-S:} Entradas: \codigo{anioNacimiento}, \codigo{anioActual}. Proceso: \texttt{edad = anioActual - anioNacimiento}; si \texttt{edad >= 18} $\rightarrow$ ``Mayor de edad'', si no $\rightarrow$ ``Menor de edad''.
  \item \textbf{Prueba:} \codigo{2008} y \codigo{2026} $\rightarrow$ \texttt{edad = 18} $\rightarrow$ ``Mayor de edad''. \codigo{2010} y \codigo{2026} $\rightarrow$ \codigo{16} $\rightarrow$ ``Menor de edad''.
\end{itemize}

\textbf{Ejercicio 7 (Error en promedio):}
\begin{itemize}
  \item \textbf{Prueba defectuosa (\texttt{a = 10, b = 20}):} Primero divide \texttt{20 / 2 = 10}, luego suma \texttt{10 + 10 = 20} (¡falso, el promedio es 15!).
  \item \textbf{Corrección:} \texttt{promedio <- (a + b) / 2.0} $\rightarrow$ \texttt{(10 + 20) / 2.0 = 15.0}.
\end{itemize}

\textbf{Ejercicio 8 (Gasolina en moto):}
\begin{itemize}
  \item \textbf{E-P-S:} Entradas: \codigo{distanciaKm}, \codigo{kmPorGalon}, \codigo{precioGalon}. Proceso: \codigo{galones = distanciaKm / kmPorGalon}, \codigo{costo = galones * precioGalon}.
  \item \textbf{Prueba (\codigo{240 km}, \codigo{120 km/gal}, \codigo{\$15500}):} \codigo{galones = 2.0}, \codigo{costo = \$31000}.
\end{itemize}
\end{solucion}

\end{document}
